% ==================================================
% 第三章　研究方法與全身模型預測控制器設計（精簡重整版）
% 圖片統一使用 PNG；本檔直接替換原第三章，不與原稿同時引入。
% ==================================================
\chapter{研究方法與全身模型預測控制器設計}
\label{c:methodology}
\label{c:Mathodology}

本章以 Tang 等人\cite{tang2024gmpc}提出之幾何模型預測控制為基礎，
結合手臂命令追蹤模型、動態障礙預測與操作參考生成，
共同求解底盤速度及手臂位置設定點變化率。
以下依序說明全身運動學、增廣預測、控制問題、
視覺導引操作及執行端整合。

% ==================================================
% 3.1 底盤幾何運動與全身運動學
% ==================================================
\section{底盤幾何運動與全身運動學}
\label{sec:wholebody_kinematics}

% ==================================================
% 3.1.1 座標系定義與位姿表示
% ==================================================
\subsection{座標系定義與位姿表示}
\label{subsec:frames_pose_representation}

定義世界、底盤、手臂基座及工具中心點座標系，
分別記為 $\mathcal F_W$、$\mathcal F_B$、$\mathcal F_A$ 與 $\mathcal F_E$。
底盤座標對應 \texttt{base\_footprint}，
$x_B$ 向前、$y_B$ 向左、$z_B$ 向上。
以 $ {}^iT_j$ 表示由座標系 $j$ 轉換至 $i$ 的齊次變換。
底盤位姿與手臂關節角為
\begin{equation}
    X_B=
    \begin{bmatrix}
        R_B(\theta_B)&p_B\\
        \mathbf 0_{1\times2}&1
    \end{bmatrix}\in\mathrm{SE}(2),
    \qquad
    p_B=\begin{bmatrix}x_B&y_B\end{bmatrix}^{\top},
    \quad q_a\in\mathbb R^6 .
    \label{eq:ch3_base_pose}
\end{equation}
其中 $R_B$ 為平面旋轉矩陣。
指數與對數映射的定義依 Sol\`a 等人\cite{sola2018micro}之整理。
如\cref{fig:ref_rotation_representation}所示，
角度主值跨越邊界時會跳變，對應旋轉仍可連續表示。

\begin{figure}[htbp]
    \centering
    \includegraphics[width=0.88\textwidth]{figsrc/ch03/rotation_representation_redrawn.png}
    \caption{旋轉與角度主值表示之比較：
    (a) 旋轉方向連續通過角度邊界；
    (b) 對應角度主值產生包絡跳變。
    依文獻\cite{tang2024gmpc}之幾何表示概念自行繪製。}
    \label{fig:ref_rotation_representation}
\end{figure}

% ==================================================
% 3.1.2 底盤 SE(2) 運動模型與追蹤誤差
% ==================================================
\subsection{底盤 SE(2) 運動模型與追蹤誤差}
\label{subsec:base_se2_model_error}

全向底盤的本體速度為
$\xi_B=[v_x^B,v_y^B,\omega_B]^{\top}$，
滿足 $\dot X_B=X_B\widehat{\xi}_B$。
令參考位姿為 $X_r$，其本體速度為 $\xi_r$，
定義幾何追蹤誤差及速度偏差為
\begin{equation}
    E_B=X_r^{-1}X_B,\qquad
    e_B=(\log E_B)^\vee,\qquad
    \delta\xi=\xi_B-\xi_r .
    \label{eq:ch3_geometric_tracking_error}
\end{equation}
其中 $(\widehat{\cdot})$ 與 $(\cdot)^\vee$ 為速度向量及李代數矩陣間的映射。
由相對變換的微分，在小誤差下建立局部模型，
並以前向歐拉法離散化：
\begin{equation}
    e_{B,k+1}=A_{B,k}e_{B,k}+\Delta t\,\delta\xi_k,
    \qquad
    A_{B,k}=I_3-\Delta t\,\operatorname{ad}(\xi_{r,k}).
    \label{eq:ch3_discrete_base_error_model}
\end{equation}
此處 $\operatorname{ad}(\xi)\eta=
([\widehat\xi,\widehat\eta])^\vee$，
$[A,B]=AB-BA$。
控制器每次更新均由目前位姿重新計算誤差，
其近似精度取決於追蹤偏差、控制週期及預測時域。

% ==================================================
% 3.1.3 移動機械臂正向運動學與全身雅可比矩陣
% ==================================================
\subsection{移動機械臂正向運動學與全身雅可比矩陣}
\label{subsec:wholebody_fk_jacobian}

將 $X_B$ 嵌入三維變換 $ {}^WT_B$，
其平移為 $[x_B,y_B,0]^{\top}$。
手臂固定安裝變換 $ {}^BT_A$ 的平移為 $(0.11,0,0.33)$ m，
已包含底盤高度與安裝偏移。
末端世界位姿為
\begin{equation}
    {}^WT_E={}^WT_B\,{}^BT_A\,{}^AT_E(q_a).
    \label{eq:ch3_wholebody_forward_kinematics}
\end{equation}
其中 $ {}^AT_E$ 包含工具中心點（Tool Center Point, TCP）偏移。
將末端線速度與角速度均表示於世界座標，可得
\begin{equation}
    {}^W\nu_E=
    \begin{bmatrix}{}^W\dot p_E\\{}^W\omega_E\end{bmatrix}
    =
    \begin{bmatrix}{}^WJ_B&{}^WJ_A\end{bmatrix}
    \begin{bmatrix}\xi_B\\\dot q_a\end{bmatrix}.
    \label{eq:ch3_wholebody_geometric_velocity}
\end{equation}
底盤雅可比將本體平移及偏航轉換為末端運動；
手臂雅可比則由各關節軸與 TCP 位置建立。
兩者共同作用的關係如\cref{fig:ch3_wholebody_motion_coupling}所示。
本研究控制輸出為設定點變化率，與實際關節速度 $\dot q_a$ 分別處理。
後續以 $\chi=[x_B,y_B,\theta_B,q_a^\top]^\top$ 表示九維全身構型。

\begin{figure}[htbp]
    \centering
    \includegraphics[width=0.98\textwidth]{figsrc/ch03/wholebody_motion_coupling.png}
    \caption{底盤側移時之全身運動學耦合：
    (a) 底盤移動前；(b) 底盤側移後，手臂調整構型以維持相同末端目標。}
    \label{fig:ch3_wholebody_motion_coupling}
\end{figure}

% ==================================================
% 3.2 手臂設定點追蹤與增廣預測模型
% ==================================================
\section{手臂設定點追蹤與增廣預測模型}
\label{sec:arm_augmented_prediction}

% ==================================================
% 3.2.1 手臂位置命令與伺服追蹤動態
% ==================================================
\subsection{手臂位置命令與伺服追蹤動態}
\label{subsec:arm_setpoint_servo}

位置介面先將設定點變化率 $u_a$ 積分為位置設定點 $s_a$，
再由關節位置控制器追蹤。
為使預測符合命令執行特性\cite{grandia2019frequency}，
以物理步長 $\Delta t_p$ 建立逐軸一階仿射模型：
\begin{equation}
    \begin{aligned}
        q_a^{j+1}&=(I_6-D_\alpha)q_a^j+D_\alpha s_a^j+b_a,\\
        s_a^{j+1}&=s_a^j+\Delta t_p u_a^j .
    \end{aligned}
    \label{eq:ch3_physics_step_arm_model}
\end{equation}
其中 $D_\alpha=\operatorname{diag}(\alpha_1,\ldots,\alpha_6)$，
$0<\alpha_i<1$，$b_a$ 為每步偏置。
$s_a^j$ 定義為作用於第 $j$ 至 $j+1$ 步的位置設定點。
設定點固定時，關節仍可能持續追蹤，
其穩態偏移為 $D_\alpha^{-1}b_a$；
因此 $u_a$ 不直接等同實際關節速度。

% ==================================================
% 3.2.2 實際關節位置與設定點之增廣狀態
% ==================================================
\subsection{實際關節位置與設定點之增廣狀態}
\label{subsec:augmented_arm_state}

令控制週期 $\Delta t=m\Delta t_p$。
在一個控制週期內保持 $u_a$ 固定，
將物理步模型合成 $m$ 次，得到
\begin{equation}
    \begin{aligned}
        q_{a,k+1}&=A_q q_{a,k}+A_s s_{a,k}+B_a u_{a,k}+c_a,\\
        s_{a,k+1}&=s_{a,k}+\Delta t u_{a,k}.
    \end{aligned}
    \label{eq:ch3_control_step_arm_model}
\end{equation}
令 $\beta_i=1-\alpha_i$，
則 $A_q=\operatorname{diag}(\beta_i^m)$、$A_s=I_6-A_q$，
$B_a$ 的第 $i$ 個對角元素為
$\alpha_i\Delta t_p\sum_{\ell=0}^{m-1}\ell\beta_i^{m-1-\ell}$，
且 $(c_a)_i=(b_a)_i(1-\beta_i^m)/\alpha_i$。

結合底盤誤差，定義
$\zeta_k=[e_{B,k}^{\top},q_{a,k}^{\top},s_{a,k}^{\top}]^{\top}
\in\mathbb R^{15}$，
$u_k=[\delta\xi_k^{\top},u_{a,k}^{\top}]^{\top}\in\mathbb R^9$。
增廣模型為
\begin{equation}
    \begin{gathered}
    \zeta_{k+1}=F_k\zeta_k+G_ku_k+d_k,\\
    F_k=
    \begin{bmatrix}
        A_{B,k}&0&0\\
        0&A_q&A_s\\
        0&0&I_6
    \end{bmatrix},
    \quad
    G_k=
    \begin{bmatrix}
        \Delta t I_3&0\\
        0&B_a\\
        0&\Delta t I_6
    \end{bmatrix}.
    \end{gathered}
    \label{eq:ch3_augmented_prediction_dynamics}
\end{equation}
其中 $d_k=[0_3^{\top},c_a^{\top},0_6^{\top}]^{\top}$。
預測末端位姿由 $X_B=X_r\exp(\widehat e_B)$
與預測的實際關節位置 $q_a$ 重建。
設定點與實際關節位置的差異如\cref{fig:ch3_arm_setpoint_tracking}所示。

\begin{figure}[htbp]
    \centering
    \includegraphics[width=0.98\textwidth]{figsrc/ch03/arm_setpoint_tracking.png}
    \caption{設定點與實際關節位置之追蹤關係：
    (a) 設定點固定；(b) 關節持續追蹤；(c) 關節逐漸接近設定點。
    虛線與實線分別表示位置設定點與實際關節位置。}
    \label{fig:ch3_arm_setpoint_tracking}
\end{figure}

% ==================================================
% 3.2.3 多步預測模型建立與參數辨識
% ==================================================
\subsection{多步預測模型建立與參數辨識}
\label{subsec:multistep_prediction_identification}

將未來 $N$ 步狀態與輸入分別堆疊為
$\mathbf Z_k=[\zeta_{1|k}^{\top},\ldots,\zeta_{N|k}^{\top}]^{\top}$
與 $\mathbf U_k=[u_{0|k}^{\top},\ldots,u_{N-1|k}^{\top}]^{\top}$，
反覆代入增廣模型可得
\begin{equation}
    \mathbf Z_k=\Phi_k\zeta_{0|k}+\Gamma_k\mathbf U_k+\eta_k .
    \label{eq:ch3_condensed_augmented_prediction}
\end{equation}
其中 $\Phi_k$、$\Gamma_k$ 與 $\eta_k$
分別描述初始狀態、控制序列及累積偏置對預測的影響。
模型參數以對齊後的設定點及關節量測逐軸進行最小平方法辨識，
並以時間分離的資料驗證單步及多步預測誤差。
延遲前推與線上偏置更新於第 3.7 節說明。

% ==================================================
% 3.3 全身預測控制之目標函數與限制條件
% ==================================================
\section{全身預測控制之目標函數與限制條件}
\label{sec:wholebody_objective_constraints}

% ==================================================
% 3.3.1 底盤與末端位姿追蹤目標
% ==================================================
\subsection{底盤與末端位姿追蹤目標}
\label{subsec:wholebody_tracking_objective}

令 $\|x\|_Q^2=x^\top Qx$。
末端目標及預測位姿均表示於世界座標，
定義六維殘差為
\begin{equation}
    r_{E,i|k}=
    \begin{bmatrix}
        p_{E,i|k}^{d}-p_{E,i|k}\\
        R_{E,i|k}
        [\log(R_{E,i|k}^{\top}R_{E,i|k}^{d})]^\vee
    \end{bmatrix}.
    \label{eq:ch3_tcp_pose_residual}
\end{equation}
姿態殘差轉換至世界座標，使其與幾何雅可比一致。
基本目標整合底盤與末端追蹤、輸入大小及命令平順性：
\begin{equation}
    \begin{aligned}
        J={}&\sum_{i=1}^{N}
        \left(\|e_{B,i|k}\|_{Q_{B,i}}^2+
              \|r_{E,i|k}\|_{W_E}^2\right)\\
        &+\sum_{i=0}^{N-1}
        \left(\|u_{i|k}\|_R^2+
              \|\Delta v_{i|k}\|_S^2\right).
    \end{aligned}
    \label{eq:ch3_wholebody_objective}
\end{equation}
其中 $Q_{B,N}$ 為終端權重，
$v=[\xi_B^\top,u_a^\top]^\top$ 為完整命令，
$\xi_B=\xi_r+\delta\xi$。
第一步差分使用預測生效起點之前的實際命令，
因此底盤命令差分包含參考速度的變化。
各權重依位置、姿態及命令的容許尺度設定；
操作階段可另加入構型引導與避障鬆弛成本。
各成本權重取半正定矩陣，輸入權重提供正則化。

% ==================================================
% 3.3.2 關節、命令變化與輪級運動限制
% ==================================================
\subsection{關節、命令變化與輪級運動限制}
\label{subsec:joint_command_wheel_constraints}

控制器同時限制實際關節位置與設定點，
並約束完整命令及相鄰命令變化。
令 $\mu\in\mathbb R_{\geq0}^{6}$ 為關節限位餘量：
\begin{equation}
    \begin{gathered}
        q_{\min}+\mu\leq q_{a,i|k}\leq q_{\max}-\mu,\qquad
        q_{\min}+\mu\leq s_{a,i|k}\leq q_{\max}-\mu,\\
        v_{\min}\leq v_{i|k}\leq v_{\max},\qquad
        -\Delta v_{\max}\leq\Delta v_{i|k}\leq\Delta v_{\max}.
    \end{gathered}
    \label{eq:ch3_joint_command_limits}
\end{equation}
對四輪十字配置，依本研究輪序及正向約定，
輪角速度與底盤速度之關係為
\begin{equation}
    \dot q_w=\frac{1}{r_w}W\xi_B,\qquad
    W=
    \begin{bmatrix}
        0&1&L\\-1&0&L\\0&-1&L\\1&0&L
    \end{bmatrix}.
    \label{eq:ch3_wheel_command_mapping}
\end{equation}
其中 $r_w$ 為輪半徑，$L$ 為車心至輪心距離。
由此施加輪速與輪速命令變化率限制：
\begin{equation}
    \left|\frac{W\xi_{B,i|k}}{r_w}\right|\leq\bar\omega_w,\qquad
    \left|\frac{W(\xi_{B,i|k}-\xi_{B,i-1|k})}{r_w}\right|
    \leq\Delta t\,\bar a_w .
    \label{eq:ch3_wheel_speed_acceleration_limits}
\end{equation}
不等式逐分量成立；兩組上限均以四維向量表示。
平移與旋轉共同占用輪速能力，
因此各軸速度框之外仍需保留輪級限制。

% ==================================================
% 3.3.3 最佳化問題建構與數值求解
% ==================================================
\subsection{最佳化問題建構與數值求解}
\label{subsec:wholebody_qp_formulation}

增廣預測為控制序列的仿射函數，
末端正向運動學則需沿目前初始控制序列建立局部近似。
將狀態預測與全身雅可比串接，可得
\begin{equation}
    r_{E,i|k}\approx r_{0,i|k}-L_{i,k}\mathbf U_k .
    \label{eq:ch3_tcp_affine_residual}
\end{equation}
其中 $r_{0,i|k}$ 包含展開點與目標的差異，
$L_{i,k}$ 描述輸入序列對末端殘差的局部影響。
設定點雖不直接出現在正向運動學，
仍透過增廣動態影響未來關節位置。
代入追蹤成本及限制後形成二次規劃：
\begin{equation}
    \begin{aligned}
        \min_{\mathbf U_k,\,\varepsilon}\quad&
        \frac12\mathbf U_k^\top H_k\mathbf U_k
        +g_k^\top\mathbf U_k+J_\varepsilon\\
        \text{subject to}\quad&
        A_{\mathrm{ineq},k}\mathbf U_k
        +B_{\varepsilon,k}\varepsilon\leq b_{\mathrm{ineq},k},
        \qquad\varepsilon\geq0 .
    \end{aligned}
    \label{eq:ch3_wholebody_qp}
\end{equation}
$\varepsilon$ 僅用於允許鬆弛的限制；
關節、速度等硬限制對應的 $B_{\varepsilon,k}$ 列為零。
目前整合控制器以 Goldfarb--Idnani 方法
\cite{goldfarb1983dual}求解，必要時使用 OSQP
\cite{stellato2020osqp}作為備援。
結果須通過數值有限性、求解狀態及限制殘差核對。

% ==================================================
% 3.3.4 滾動時域控制原則
% ==================================================
\subsection{滾動時域控制原則}
\label{subsec:wholebody_receding_horizon}

每次更新依新狀態、參考及障礙資訊重建控制問題，
並將上一輪可沿用的控制序列向前移作為初始解。
候選結果的首步底盤命令為
$\xi_{B,k}^{\mathrm{cmd}}=\xi_{r,k}+\delta\xi_{0|k}^{*}$，
手臂命令為 $u_{a,0|k}^{*}$。
在非同步實作中，控制段須通過接受檢查並於指定生效步接續；
求解完成與命令實際生效分別記錄。

% ==================================================
% 3.4 動態障礙預測與導航參考調整
% ==================================================
\section{動態障礙預測與導航參考調整}
\label{sec:dynamic_obstacle_avoidance}

% ==================================================
% 3.4.1 障礙物表示與運動估測
% ==================================================
\subsection{障礙物表示與運動估測}
\label{subsec:obstacle_representation_prediction}

雷射回波經有效距離、自身回波篩選及座標轉換後形成群集，
再與既有軌跡進行關聯。
以群集中心估計運動，並保留物件表面樣本作為避障幾何。
對第 $g$ 個物件，採用等速度卡爾曼模型
\cite{kalman1960filter}：
\begin{equation}
    x_g^+=
    \begin{bmatrix}I_2&\Delta\tau I_2\\0&I_2\end{bmatrix}x_g+w_g,
    \qquad
    y_g=\begin{bmatrix}I_2&0\end{bmatrix}x_g+n_g .
    \label{eq:ch3_obstacle_tracking_model}
\end{equation}
其中 $x_g=[c_g^\top,v_g^\top]^\top$，
$\Delta\tau$ 為追蹤更新間隔，$w_g,n_g$ 為雜訊。
預測與量測修正的循環如\cref{fig:ref_kalman_predict_update}所示。
同一物件的表面樣本共用估計平移速度；
短暫失去觀測時依中心位移更新樣本，並限制保留時間。
此近似未描述遮蔽期間的物件旋轉與變形。

\begin{figure}[htbp]
    \centering
    \includegraphics[width=0.98\textwidth]{figsrc/ch03/kalman_tracking_redrawn.png}
    \caption{障礙追蹤之卡爾曼預測與量測更新循環，
    依文獻\cite{welch2004kalman}之方法自行繪製。}
    \label{fig:ref_kalman_predict_update}
\end{figure}

參考調整層另外由可見中心的歷史位移估計速度，
並以一階濾波平滑；歷史不足時由追蹤器初始化。
該速度用於候選運動評估，CBF 則使用追蹤器的樣本速度，
兩者的資料來源與更新時序分別保留。

% ==================================================
% 3.4.2 未來相對運動與衝突預測
% ==================================================
\subsection{未來相對運動與衝突預測}
\label{subsec:future_motion_conflict}

對候選 $c$ 生成底盤位置 $p_{B,i}^{(c)}$，
並將各障礙群依其速度外推至相同未來時間。
以 $o_{\ell,i}$ 表示群內保留的表面樣本，
定義整段預測的最小間距餘量：
\begin{equation}
    o_{\ell,i}=\hat o_{\ell,k}+\tau_i\tilde v_{g(\ell),k},
    \qquad
    D_c=\min_{\ell,i}
    \left(\|p_{B,i}^{(c)}-o_{\ell,i}\|-\rho_\ell\right).
    \label{eq:ch3_candidate_clearance_margin}
\end{equation}
$\rho_\ell$ 包含底盤外形與裕度。
使用表面樣本時不重複加入完整物件半徑；
中心表示則加入對應外形半徑。
比較涵蓋同一候選視窗內的相關障礙，
如\cref{fig:ch3_dynamic_obstacle_prediction}所示。

\begin{figure}[htbp]
    \centering
    \includegraphics[width=0.98\textwidth]{figsrc/ch03/dynamic_obstacle_prediction.png}
    \caption{動態障礙之候選運動比較：
    (a) 目前觀測；(b) 維持原速度；(c) 降低速度；(d) 側向避讓。
    (b)--(d) 為同一未來時刻；虛線圓表示目前底盤位置。}
    \label{fig:ch3_dynamic_obstacle_prediction}
\end{figure}

% ==================================================
% 3.4.3 速度與側移候選評估及導航參考更新
% ==================================================
\subsection{速度與側移候選評估及導航參考更新}
\label{subsec:candidate_reference_update}

候選組合包含速度縮放、側向偏移及路徑推進調整。
對通過間距檢查的候選，
以路徑進度、偏離量及切換成本選擇；
若均未通過，則選擇間距較大的退化候選，
配合減速或停止處理。
候選切換保留遲滯，以降低感測波動造成的反覆變更。
此層提供參考調整，實際命令仍受 QP 與執行端限制。

位置及朝向參考更新後，由相鄰位姿重建參考速度：
\begin{equation}
    \xi_{r,i}=
    \begin{bmatrix}
        R_B(\theta_{r,i})^\top
        (p_{r,i+1}-p_{r,i})/\Delta t\\
        \operatorname{wrap}(\theta_{r,i+1}-\theta_{r,i})/\Delta t
    \end{bmatrix}.
    \label{eq:ch3_reference_velocity_reconstruction}
\end{equation}
其中 $\operatorname{wrap}$ 將角度差映射至 $(-\pi,\pi]$。
此有限差分近似與離散參考採用相同時間基準。

% ==================================================
% 3.4.4 控制屏障函數與避障限制
% ==================================================
\subsection{控制屏障函數與避障限制}
\label{subsec:cbf_obstacle_constraints}

控制屏障函數（Control Barrier Function, CBF）
用於將障礙間距轉換為控制限制\cite{ames2019cbf}。
對樣本 $\ell$ 定義
$h_\ell=\|p_B-o_\ell\|^2-\rho_\ell^2$。
在目前預測幾何上固定相對向量
$\bar d_{\ell,i}=\bar p_{B,i}-o_{\ell,i}$，
可建立仿射速度限制：
\begin{equation}
    2\bar d_{\ell,i}^{\top}
    \left(R_B(\bar\theta_{B,i})
    \begin{bmatrix}v_x^B\\v_y^B\end{bmatrix}_{i|k}
    -\hat v_{\ell,k}\right)
    +\alpha_h\bar h_{\ell,i}
    \geq-\varepsilon_{c(\ell),i|k},
    \qquad \varepsilon_{c(\ell),i|k}\geq0 .
    \label{eq:ch3_cbf_affine_constraint}
\end{equation}
其中 $\alpha_h>0$，$c(\ell)$ 表示樣本所屬的限制群組；
同群組樣本依配置共用鬆弛，
各群組與預測步的鬆弛以加權平方成本加入目標函數。
限制中的底盤速度為完整速度 $\xi_r+\delta\xi$。
實作採局部幾何、離散更新及有限裕度，
且可能啟用鬆弛，因此仍需核對預測間距、鬆弛量及實際軌跡。

% ==================================================
% 3.4.5 原始雷射掃描防護與命令處理
% ==================================================
\subsection{原始雷射掃描防護與命令處理}
\label{subsec:raw_scan_command_protection}

執行路徑另以原始有效回波建立近距離防護，
處理未進入追蹤器的障礙。
令 $d_\ell$ 為回波至圓盤底盤表面的間距，
$n_\ell$ 為朝向回波的單位向量，
$v_{\mathrm{app},\ell}$ 為輸入平移速度沿該方向的非負投影。
保護距離及輸出限制為
\begin{equation}
    \begin{aligned}
        d_{\mathrm{stop},\ell}
        &=d_0+v_{\mathrm{app},\ell}(\tau_0+t_{\mathrm{age}})
          +\frac{v_{\mathrm{app},\ell}^2}{2a_{\mathrm{brake}}}
          +\epsilon_{\mathrm{scan}}+\epsilon_{\mathrm{body}},\\
        n_\ell^\top
        \begin{bmatrix}v_x^{B,\mathrm{out}}\\v_y^{B,\mathrm{out}}\end{bmatrix}
        &\leq\alpha_{\mathrm{sh}}(d_\ell-d_{\mathrm{stop},\ell}).
    \end{aligned}
    \label{eq:ch3_scan_stopping_distance}
\end{equation}
各參數分別描述基本間距、延遲與掃描年齡、減速能力及近似裕度。
圓盤模型的法向旋轉速度為零，未直接涵蓋伸展手臂的掃掠範圍。

防護命令以有限次半空間投影處理，完成後檢查殘差。
若需放寬主動退離要求為非接近條件，另外記錄該退化處理；
仍不符合條件時請求停止。
掃描逾時則限制速度。
命令依序通過平滑、防護與輪級限制，
最終放行值回報予預測模型，
其作用位置如\cref{fig:ch3_navigation_command_layers}所示。
操作門扇的回波例外與專用機構幾何檢查共同配置。

\begin{figure}[htbp]
    \centering
    \includegraphics[width=0.98\textwidth]{figsrc/ch03/navigation_command_layers.png}
    \caption{導航參考調整、QP 避障限制與執行端防護之作用位置。}
    \label{fig:ch3_navigation_command_layers}
\end{figure}

% ==================================================
% 3.5 視覺目標估測與任務參考生成
% ==================================================
\section{視覺目標估測與任務參考生成}
\label{sec:visual_task_reference}

% ==================================================
% 3.5.1 任務目標之視覺定位與幾何估測
% ==================================================
\subsection{任務目標之視覺定位與幾何估測}
\label{subsec:visual_geometric_estimation}

本研究以 Mask R-CNN\cite{he2017maskrcnn}
分割門把、門扇、圓柱、抽屜把手、筆狀代理物及抽屜面板六類目標。
實作採用 ResNet-50-FPN v2，
以 COCO 預訓練權重初始化，
再以 Isaac Sim 的 RGB 影像及語意標註建立實例資料。
線上推論僅輸入 RGB；
信心分數與遮罩二值化門檻皆為 $0.5$。
資料配置及 IoU 評估留於第二章與第四章說明，
目前模型的訓練與測試範圍皆為模擬影像。

完整架構如\cref{fig:ref_mask_rcnn_framework}所示，
並依 Torchvision 的 v2 實作核對\cite{torchvisionmaskrcnnv2}。
ResNet-50\cite{he2016resnet}的四個殘差階段輸出 C2--C5，
FPN\cite{lin2017fpn}建立多尺度特徵；
RPN 產生候選框後，偵測分支進行類別與外框迴歸，
遮罩分支則以 RoIAlign 特徵產生各類別的像素遮罩。
本研究的預測層包含六個前景類別與背景，共七類。
圖中 Conv 區塊包含對應的正規化與啟用操作；
網路作為感知前端，控制方法承接其分割結果。

\begin{figure}[htbp]
    \centering
    \includegraphics[width=0.98\textwidth]{figsrc/ch03/mask_rcnn_architecture.png}
    \caption{本研究使用之 Mask R-CNN ResNet-50-FPN v2 推論架構：
    (a) CNN 骨幹與特徵金字塔；(b) 區域候選網路；
    (c) 類別與外框分支；(d) 遮罩分支。
    依文獻\cite{he2017maskrcnn}及本研究實作自行繪製；
    遮罩分支於推論時使用偵測後的外框。}
    \label{fig:ref_mask_rcnn_framework}
\end{figure}

實例遮罩用於選取有效深度點，
門把則由遮罩外框建立擴張 ROI，以保留門板與橫桿周圍資訊。
令 $K_C$ 為相機內參、$z(u,v)$ 為光軸深度，
三維點由 $ {}^Cp=z(u,v)K_C^{-1}[u,v,1]^\top$ 反投影取得。
各類物件依幾何先驗估測：
圓柱以已知半徑進行平面圓擬合；
筆狀物件由主成分取得長軸；
抽屜把手結合面板方向與橫桿點群；
門扇先依鉸鏈徑向、角度與高度篩選點群，
再由主要角度區域估計開度。
尺寸、支持點跨度及幾何一致性用於排除不合理結果。

% ==================================================
% 3.5.2 影像、機器人狀態與座標轉換之時間對齊
% ==================================================
\subsection{影像、機器人狀態與座標轉換之時間對齊}
\label{subsec:visual_state_time_alignment}

機器人在影像取得與控制求解間仍可能移動，
因此使用觀測時刻 $t_o$ 的相機位姿轉換目標：
\begin{equation}
    \begin{bmatrix}{}^Wp\\1\end{bmatrix}
    ={}^WT_B(t_o)\,{}^BT_C(q_a(t_o))
    \begin{bmatrix}{}^Cp\\1\end{bmatrix}.
    \label{eq:ch3_visual_timestamp_transform}
\end{equation}
固定相機使用固定外參；手腕相機的變換包含當時手臂構型與安裝偏移。
目標封包保留原觀測時間、座標系、幾何量及有效性，
接近前核對資料年齡、可見點數與幾何一致性。
短暫遮蔽僅在設定期限內保留估計，
超過條件則重新觀測或停止推進。

% ==================================================
% 3.5.3 可行操作區域與底盤目標選擇
% ==================================================
\subsection{可行操作區域與底盤目標選擇}
\label{subsec:feasible_operation_base_selection}

依目標及家具幾何生成有限候選，
檢查導航可達性、底盤淨空、關節限制及預抓與抓取構型。
移動中操作另沿候選區域取樣，
以前一關節解初始化下一點，核對構型連續性與全身淨空。
對通過檢查的集合 $\mathcal F_k$，選位成本為
\begin{equation}
    \begin{aligned}
        (b^*,q_a^*)&=\underset{(b,q_a)\in\mathcal F_k}{\operatorname{arg\,min}}
        J_{\mathrm{sel}}(b,q_a),\\
        J_{\mathrm{sel}}&=
        w_LL_{\mathrm{route}}
        +\frac{w_\sigma}{\max(\sigma_{\mathrm{path}},\epsilon_\sigma)}
        +w_{\mathrm{ext}}[\ell_{\mathrm{ext}}-\ell_{\mathrm{comf}}]_+
        +J_{\mathrm{heading}} .
    \end{aligned}
    \label{eq:ch3_operation_base_selection}
\end{equation}
$L_{\mathrm{route}}$ 包含進入操作區及離開後前往下一目標的路程，
以 Dijkstra 方法\cite{dijkstra1959note}估計；
$\sigma_{\mathrm{path}}$ 為取樣構型之選位雅可比最小奇異值的最小值。
$\ell_{\mathrm{ext}}$ 為手臂基座至目標的水平距離，
超過舒適距離 $\ell_{\mathrm{comf}}$ 才增加伸展成本，
$[x]_+=\max(x,0)$。
$J_{\mathrm{heading}}$ 考量進入與後續行進方向的角度差。
為減少逆運動學計算，候選依成本下界排序並剪枝。
選擇範圍限於已生成且通過檢查的候選，
實際接近路徑仍由導航前端更新。

\begin{figure}[htbp]
    \centering
    \includegraphics[width=0.98\textwidth]{figsrc/ch03/visual_operation_reference.png}
    \caption{操作目標導引之底盤選位：
    (a) 選位前的機器人與目標位置；(b) 選位後，手臂接近目標。}
    \label{fig:ch3_visual_operation_reference}
\end{figure}

% ==================================================
% 3.5.4 導航路徑下之朝向與末端參考生成
% ==================================================
\subsection{導航路徑下之朝向與末端參考生成}
\label{subsec:navigation_manipulation_reference}

接近階段沿用導航路徑，
再依選定操作位姿調整局部側移、朝向與速度。
令 $\hat n_i$ 為路徑法向，
$\ell_{\perp,i}$ 為操作底盤位置相對路徑的法向差，
位置與朝向參考寫為
\begin{equation}
    \begin{aligned}
        p_{r,i}&=p_{\mathrm{nav},i}
        +a_k w_{\mathrm{pull},i}
        \operatorname{clip}(\ell_{\perp,i},-\ell_{\max},\ell_{\max})\hat n_i,\\
        \theta_{r,i}&=\theta_{\mathrm{seed},i}
        +w_i\operatorname{wrap}(\theta_g-\theta_{\mathrm{seed},i}).
    \end{aligned}
    \label{eq:ch3_navigation_local_pull}
\end{equation}
側移距離權重 $w_{\mathrm{pull},i}$ 依切向距離漸增，
時間權重 $a_k$ 以三次函數 $S(x)=3x^2-2x^3$ 平順接手。
$\theta_{\mathrm{seed},i}$ 承接操作開始時的朝向，
$w_i$ 依距離漸增至操作朝向 $\theta_g$。
轉入收納導航時，保留切換當下的側移向量並漸減至零，
避免路徑方向更新造成退回方向跳變。
參考依合成平移速度及偏航速率重取樣，
再由式\eqref{eq:ch3_reference_velocity_reconstruction}重建速度。

手臂以前段關節空間展開、後段併入世界抓取線。
定義五次平滑函數
$\psi(\nu)=10\nu^3-15\nu^4+6\nu^5$
\cite{flash1985coordination}，
由距離與時間共同限制進度：
\begin{equation}
    \nu_i=
    \min\left\{
        \operatorname{clip}\left(
        \frac{d_{\mathrm{out}}-d_i}{d_{\mathrm{out}}-d_{\mathrm{in}}},0,1\right),
        \nu_k+\frac{t_i-t_k}{T_{\mathrm{exp,min}}}
    \right\},
    \qquad \beta_i=\psi(\nu_i).
    \label{eq:ch3_distance_based_arm_progress}
\end{equation}
$d_i$ 為底盤參考至操作位置的距離，$\nu_k$ 為承接進度。
最短展開時間取
$T_{\mathrm{exp,min}}=\max(T_0,
1.875\|q_a^g-q_{a,0}\|_\infty/\bar v_{q,\mathrm{target}})$，
避免大幅關節位移在過短時間內展開。
實際進度受距離、時間及狀態共同更新，
不將整段運動視為固定時間的最小急動軌跡。

由起始構型 $q_{a,0}$ 與選位所得構型 $q_a^g$ 建立
\begin{equation}
    q_{a,i}^{\mathrm{exp}}=q_{a,0}+\beta_i(q_a^g-q_{a,0}),\qquad
    T_{E,i}^{\mathrm{exp}}=
    \operatorname{FK}_{\mathrm{WB}}(b_{r,i},q_{a,i}^{\mathrm{exp}}).
    \label{eq:ch3_joint_space_unfold_reference}
\end{equation}
$b_{r,i}$ 為同時刻的底盤參考位姿，
$\operatorname{FK}_{\mathrm{WB}}$ 依式\eqref{eq:ch3_wholebody_forward_kinematics}計算。
啟用構型引導時，另加入
$J_{\mathrm{post}}=\sum_i
\|q_{a,i|k}-q_{a,i}^{\mathrm{exp}}\|_{Q_{\mathrm{post}}}^2$。
展開後段以
$\gamma_i=\psi(\operatorname{clip}((\beta_i-\beta_0)/(1-\beta_0),0,1))$
併入抓取線：
\begin{equation}
    \begin{aligned}
        p_{E,i}^{d}&=(1-\gamma_i)p_{\mathrm{exp},i}
                     +\gamma_i p_{\mathrm{line},i},\\
        R_{E,i}^{d}&=R_{\mathrm{exp},i}
        \exp\left(\gamma_i\log(R_{\mathrm{exp},i}^{\top}R_{\mathrm{line},i})\right).
    \end{aligned}
    \label{eq:ch3_continuous_approach_reference}
\end{equation}
抓取線依底盤接近程度與 TCP 到達預抓位置的狀態推進，
後續操作銜接於第 3.6 節說明。
距離規則如\cref{fig:ch3_mobile_approach_distance_profile}所示；
參考平順性之外，實際命令仍受關節與輪級限制。

\begin{figure}[htbp]
    \centering
    \includegraphics[width=0.98\textwidth]{figsrc/ch03/mobile_approach_distance_profile.png}
    \caption{接近階段之距離規則示例：
    (a) 平移逐軸速度上限；(b) 手臂距離展開進度與後段併線門檻。
    圖為設計曲線，實際進度另受時間與狀態限制。}
    \label{fig:ch3_mobile_approach_distance_profile}
\end{figure}

% ==================================================
% 3.6 移動中操作與機構約束控制
% ==================================================
\section{移動中操作與機構約束控制}
\label{sec:mobile_constrained_operation}

本節以移動中操作\cite{burgess2023motm}
及關節式物件操作\cite{mittal2022articulated}為背景，
說明共用全身控制器如何承接不同任務參考。

% ==================================================
% 3.6.1 移動中接近與操作銜接
% ==================================================
\subsection{移動中接近與操作銜接}
\label{subsec:moving_approach_grasp}

接近目標以世界座標表示，
手臂配合底盤運動補償相對位移。
末端到達操作位置後，任務流程銜接夾爪動作、抬升及後續操作。
本文假設抓持後物件與末端的相對位姿保持固定，
以建立持物與機構操作的運動關係。
移動中接近仍受可達範圍與操作時間窗限制；
未能完成時依任務配置重新接近、轉為定點操作或停止。
各階段以末端到位、操作進度及執行狀態銜接，
實際抓持狀態辨識與滑移處理列為未來工作。

% ==================================================
% 3.6.2 門與抽屜之機構約束建模
% ==================================================
\subsection{門與抽屜之機構約束建模}
\label{subsec:articulated_object_constraints}

令 $\eta$ 為物件機構座標、$ {}^WT_H(\eta)$ 為把手位姿。
操作開始時，由當下構型與物件估計建立末端相對把手的變換：
\begin{equation}
    {}^HT_E=({}^WT_H(\eta_b))^{-1}{}^WT_E(\chi_b),
    \qquad
    {}^WT_E(\chi)={}^WT_H(\eta)\,{}^HT_E .
    \label{eq:ch3_grasp_closed_chain}
\end{equation}
其中 $\chi_b,\eta_b$ 為建立該關係時的狀態。
底盤與手臂需共同維持此相容關係。
對繞鉸鏈點 $h$、單位軸 $a_h$ 旋轉的門扇，
把手模型為
\begin{equation}
    \begin{aligned}
        p_H(\phi)&=h+R_{a_h}(\phi-\phi_0)(p_{H,0}-h),\\
        R_H(\phi)&=R_{a_h}(\phi-\phi_0)R_{H,0}.
    \end{aligned}
    \label{eq:ch3_revolute_handle_model}
\end{equation}
軸方向依正門角定義，門角參考經此模型轉換為末端目標。
開門時可在選定預測步加入六維位姿容許範圍：
\begin{equation}
    -\epsilon_c\leq r_{0,i|k}-L_{i,k}\mathbf U_k\leq\epsilon_c,
    \qquad i\in\mathcal I_c .
    \label{eq:ch3_contact_pose_tube}
\end{equation}
其中 $\epsilon_c$ 分別設定位置與姿態容許量，
並與底盤及手臂限制共同納入 QP。

\begin{figure}[htbp]
    \centering
    \includegraphics[width=0.98\textwidth]{figsrc/ch03/door_operation_sequence.png}
    \caption{門扇操作之全身構型：
    (a) 操作起始門角；(b) 門扇開啟後，底盤與手臂配合把手運動。}
    \label{fig:ch3_door_operation_sequence}
\end{figure}

回程再次接近前，利用視覺估計門角並重建把手目標。
關門階段採用繞鉸鏈的共轉參考：
\begin{equation}
    {}^WT_{B,r}(t)=
    \mathcal T_h(\phi_r(t)-\phi_b)\,{}^WT_{B,b},
    \qquad {}^BT_E^d={}^BT_{E,b}.
    \label{eq:ch3_door_corotation_reference}
\end{equation}
$\mathcal T_h$ 為繞鉸鏈的世界座標剛體變換。
底盤移動配合旋轉，手臂維持相對底盤的末端構型；
操作時長依輪速與變化限制設定。
停止檢查另核對 TCP 至鉸鏈的水平半徑偏差及環境淨空。
模擬執行使用回報門角核對進度與幾何，
回程再接近目標則由視覺估計建立，兩者資料用途分列。

對沿單位方向 $a_d$ 平移的抽屜，
位移參考 $d$ 對應
\begin{equation}
    p_H(d)=p_{H,0}+d a_d,\qquad R_H(d)=R_{H,0}.
    \label{eq:ch3_prismatic_handle_model}
\end{equation}
拉開與推回以位移參考方向切換，
並由相對起始 TCP 的軸向位移估計操作進度。
取筆收納後重新觀測把手，再銜接推回參考；
終點保留行程餘量，解除操作後的殘餘開度另行評估。
各階段分別配置參考、QP 限制與執行端幾何核對，
不視為全程啟用相同硬限制。

\begin{figure}[htbp]
    \centering
    \includegraphics[width=0.98\textwidth]{figsrc/ch03/drawer_operation_sequence.png}
    \caption{抽屜操作之全身構型：
    (a) 抽屜初始位置；(b) 沿滑軌拉出後，底盤與手臂配合調整。}
    \label{fig:ch3_drawer_operation_sequence}
\end{figure}

% ==================================================
% 3.6.3 持物運送與收納之運動參考
% ==================================================
\subsection{持物運送與收納之運動參考}
\label{subsec:carrying_storage_reference}

進入持物階段後，以底盤相對目標描述抬升及保持：
\begin{equation}
    {}^WT_E^d={}^WT_B\,{}^BT_E^d,\qquad
    {}^BT_E(q_a)={}^BT_A\,{}^AT_E(q_a).
    \label{eq:ch3_base_relative_carry_target}
\end{equation}
追蹤殘差直接於底盤座標計算；
目標與末端隨底盤共同運動，
因此其增廣幾何映射為
$\mathcal J_E^{\mathrm{rel}}=[0_{6\times3},{}^BJ_A,0_{6\times6}]$。
抬升、退出家具及收納以分階段構型參考銜接，
關節內插採用前述五次平滑函數，
再由正向運動學產生相對末端目標。
導航可依階段配置同時執行。

\begin{figure}[htbp]
    \centering
    \includegraphics[width=0.98\textwidth]{figsrc/ch03/base_relative_carrying.png}
    \caption{相對底盤目標下之持物運送：
    (a) 底盤移動前；(b) 底盤平移後，維持相同的底盤相對末端構型。}
    \label{fig:ch3_base_relative_carrying}
\end{figure}

到達收納構型後依序釋放物件、等待落下及收臂。
流程中的釋放事件與物件實際落入盒內分別記錄。
線上幾何檢查涵蓋啟用的機器人、家具及收納盒模型；
持物本體與盒內既有物件尚未完整納入所有階段的線上碰撞模型，
其影響於第四章說明。

% ==================================================
% 3.6.4 任務階段切換與控制狀態延續
% ==================================================
\subsection{任務階段切換與控制狀態延續}
\label{subsec:task_transition_continuity}

住宅整合任務依序執行開門通行、
茶几圓柱移動中抓取與車尾收納、
開抽屜取筆與收納、關抽屜、
回程推門通行與關門，最後返回起點。
任務層依觀測有效性、位姿誤差、操作進度及執行狀態，
更新參考、權重與啟用限制。

切換時承接實際關節位置、設定點、在途命令、
最近放行輸入及適用的模型偏置；
朝向、局部側移及展開進度依新階段延續或重新初始化。
新的操作參考由切換當下的構型建立。
若尚未取得有效目標或未通過交接檢查，
則沿目前已接受的控制段與停止備援處理。

% ==================================================
% 3.7 執行動態補償與控制實現
% ==================================================
\section{執行動態補償與控制實現}
\label{sec:execution_compensation}

% ==================================================
% 3.7.1 在途命令與實際放行輸入之預測整合
% ==================================================
\subsection{在途命令與實際放行輸入之預測整合}
\label{subsec:queued_applied_command_prediction}

新控制段生效前，既定命令仍持續作用。
令 $z^{\mathrm{ex}}=[\chi^\top,s_a^\top]^\top$
為底盤位姿、實際關節位置及設定點，
$D_k$ 為生效前已排定的物理步數。
以物理模型 $f_p$ 依命令順序前推：
\begin{equation}
    z_{0|k}^{\mathrm{ex}}=z_k^{\mathrm{ex}},\qquad
    z_{\ell+1|k}^{\mathrm{ex}}=
    f_p(z_{\ell|k}^{\mathrm{ex}},v_{\ell|k}^{\mathrm{fix}}),
    \quad \ell=0,\ldots,D_k-1 .
    \label{eq:ch3_queued_command_forward_prediction}
\end{equation}
模型結合底盤運動與設定點追蹤，
並納入已排定的命令斜坡及限制處理。
由前推終點的底盤位姿、實際關節位置及設定點建立 $\zeta_{0|k}$；
底盤誤差依式\eqref{eq:ch3_geometric_tracking_error}，
相對 $X_r(t_k^{\mathrm{eff}})$ 計算，
其中 $t_k^{\mathrm{eff}}=t_k+D_k\Delta t_p$。
參考亦由此時刻起算。
未來尚未知的防護修正則在實際放行後，
利用回報命令、設定點及新狀態更新後續預測。

% ==================================================
% 3.7.2 模型偏置估測與末端追蹤補償
% ==================================================
\subsection{模型偏置估測與末端追蹤補償}
\label{subsec:model_bias_tcp_compensation}

將每步偏置表示為 $b_a=D_\alpha d_a$，
$d_a$ 為等效關節偏移。
令一次估測區間包含 $m_o$ 個物理步，
$B_\ell=\operatorname{diag}(\beta_i^\ell)$，
$\Lambda_{m_o}=I_6-B_{m_o}$。
使用更新前的估計 $\hat d_a^-$ 與實際作用的設定點，
預測區間終點：
\begin{equation}
    \hat q_a^{j+m_o,-}=
    B_{m_o}q_a^j+
    \sum_{\ell=0}^{m_o-1}
    B_{m_o-1-\ell}D_\alpha\tilde s_a^{j+\ell}
    +\Lambda_{m_o}\hat d_a^- .
    \label{eq:ch3_interval_offset_prediction}
\end{equation}
由終點殘差取得等效偏移觀測，再平滑投影：
\begin{equation}
    \begin{aligned}
        d_a^{\mathrm{obs}}&=\hat d_a^-+
        \Lambda_{m_o}^{-1}(q_a^{j+m_o}-\hat q_a^{j+m_o,-}),\\
        \hat d_a^+&=\Pi_{\mathcal D}
        [\hat d_a^-+\lambda_d(d_a^{\mathrm{obs}}-\hat d_a^-)],
        \qquad \lambda_d=\frac{m_o\Delta t_p}{\tau_d}\in(0,1].
    \end{aligned}
    \label{eq:ch3_interval_offset_update}
\end{equation}
$\Pi_{\mathcal D}$ 同時限制絕對偏移與相對初始值的修正量。
僅使用時間對齊、數值有限且未重複使用的完整區間。
偏置於啟用的接近階段更新，夾爪動作開始後依配置凍結；
同次控制的 QP、在途命令前推及停止檢查使用一致參數。

接近末段另可對 TCP 位置目標施加有界積分修正：
\begin{equation}
    \iota_k^+=\operatorname{clip}
    (\iota_k^-+K_I^d(p_{E,k}^d-p_{E,k}),-\bar\iota,\bar\iota),
    \qquad
    p_{E,k}^{d,\mathrm{cmd}}=p_{E,k}^d+\iota_k^+ .
    \label{eq:ch3_bounded_tcp_target_compensation}
\end{equation}
修正僅在接近進度、誤差及低 TCP 速度等條件符合時更新，
並保留狀態來源與時間。
到位條件仍依原始幾何目標判定。
模型偏置與目標補償分別改變狀態轉移及追蹤目標，
兩者的啟用方式與效果於第四章分開評估。

% ==================================================
% 3.7.3 命令放行、停止軌跡檢查與逾時處置
% ==================================================
\subsection{命令放行、停止軌跡檢查與逾時處置}
\label{subsec:command_acceptance_stopping}

求解端以狀態快照產生帶版本、生效步及停止尾段的候選封包，
平台在計算期間持續執行已接受控制段。
新封包須在期限內到達，並通過版本、數值、
幾何及交接檢查，才於指定步接續。
停止軌跡保留在途命令、底盤斜坡、輪級限制及設定點追蹤動態，
核對實際關節位置、設定點、啟用幾何的淨空，
以及閉鏈位姿或共轉半徑等機構相容條件。
幾何涵蓋範圍沿用第 3.6 節的設定與限制。

晚到、失敗或未通過檢查的候選予以拒絕，
保留仍有效的排程，必要時執行已接受的停止尾段；
無法承接時啟動中止處理。
設定點變化率歸零後，關節仍可能繼續追蹤，
因此命令歸零與實際停穩分別記錄。
模擬時間與牆鐘計算耗時亦分開統計。

% ==================================================
% 3.7.4 整體控制流程與演算法
% ==================================================
\subsection{整體控制流程與演算法}
\label{subsec:overall_control_algorithm}

整體閉迴路如\cref{fig:ch3_wholebody_control_workflow}所示。
任務層更新參考與限制，共用模型提供全身預測，
執行端則依接受條件接續控制段並回報最終輸入。
演算法\ref{alg:wholebody_predictive_control}
整理背景計算與每一物理步的執行關係。

\begin{figure}[htbp]
    \centering
    \includegraphics[width=0.98\textwidth]{figsrc/ch03/wholebody_control_workflow.png}
    \caption{全身預測控制之計算、命令核對與執行回授流程。}
    \label{fig:ch3_wholebody_control_workflow}
\end{figure}

\begin{algorithm}[htbp]
    \small
    \setstretch{1.1}
    \caption{具命令承接與停止備援之全身預測控制}
    \label{alg:wholebody_predictive_control}
    \begin{algorithmic}[1]
        \Require State and sensor data, task state, navigation path,
                 scheduled commands, and model parameters
        \Ensure Applied commands, setpoints, task events, and stopping records
        \State Preserve the accepted segment and its stopping tail
        \While{the task is active or stopping is incomplete}
            \If{an update is due and the solver is available}
                \State Validate data and update enabled bias estimates
                \State Propagate scheduled commands to the activation time
                \State Update targets, references, and obstacle predictions
                \State Launch asynchronous QP solving and trajectory checks
            \EndIf
            \If{a candidate is ready and all checks pass}
                \State Submit a versioned packet with its activation step and stopping tail
            \EndIf
            \Statex \textbf{At every physics step:}
            \If{a valid packet reaches its activation step}
                \State Activate the new segment and its stopping tail
            \ElsIf{stopping is requested or the segment is nearly exhausted}
                \State Execute the applicable stopping tail; abort if necessary
            \EndIf
            \State Process base commands and integrate arm setpoints
            \State Report applied inputs, states, and task events
        \EndWhile
    \end{algorithmic}
\end{algorithm}

第四章依模型預測、動態導航、協同操作及住宅整合任務，
說明實際配置、比較條件與驗證結果。
